% Distillation with authenticated noise, then potentially corrupted codewords.
\begin{frame}{Distillation Challenges}
  \centering
  \setbeamercolor{itemize item}{fg=CVQCTeal}
  \begin{tikzpicture}[>={Latex[length=2.8mm]}]
    \path[use as bounding box] (-6.6,-3.15) rectangle (6.6,3.05);

    \node[anchor=north west,inner sep=0pt,text=CVQCInk] at (-6.35,2.45) {%
      \begin{minipage}{6.8cm}
        \raggedright
        \begin{itemize}
          \setlength{\itemsep}{1.0em}
          \item<1-> MPQC
            \href{https://doi.org/10.1007/978-3-030-45727-3_25}{[DGJ+20]}:
            modular distillation.
          \item<2-> MBQC does not guarantee correct QAS.
          \item<3-> $\Rightarrow$ Adversary can cheat during distillation.
        \end{itemize}
      \end{minipage}%
    };

    \begin{scope}[shift={(3.85,1.35)},scale=0.75,transform shape,
        magic/.style={font=\small,text=CVQCInk},
        flow/.style={->,very thick,draw=CVQCTeal}]
      \node[magic] at (-2.65,0.90)
        {$\mathsf{TEnc}(\ket{\mathrm{CCZ}})$};
      \node[magic] at (-0.45,0.90)
        {$\mathsf{TEnc}(\ket{\mathrm{CCZ}})$};
      \node[font=\large,text=CVQCMuted] at (1.15,0.90) {$\cdots$};
      \node[magic] (oldmagic) at (2.65,0.90)
        {$\mathsf{TEnc}(\textcolor{CVQCCoral}{\ket{\mathrm{CCZ}}})$};
      \only<2->{
        % Replace valid authentication of a noisy state with a bad codeword.
        \draw[draw=CVQCCoral,line width=0.9pt,line cap=round]
          (oldmagic.south west) -- (oldmagic.north east)
          (oldmagic.north west) -- (oldmagic.south east);
        \node[magic,text=CVQCCoral,above=0.25cm of oldmagic]
          {$\mathsf{TEnc}(\ket{\mathrm{CCZ}})$};
      }

      \draw[flow] (-2.65,0.60) -- (-2.20,0.15);
      \draw[flow] (-0.45,0.60) -- (-0.45,0.15);
      \draw[flow,draw=CVQCCoral] (2.65,0.60) -- (2.20,0.15);

      \path[draw=CVQCInk,fill=CVQCInk,very thick,rounded corners=3pt]
        (-3.00,0.15) -- (3.00,0.15) -- (1.10,-1.20) -- (-1.10,-1.20) -- cycle;
      \node[font=\large\bfseries,text=white] at (0,-0.50)
        {\only<3>{\color{CVQCCoral!55!white}}Distillation};
      \draw[flow] (0,-1.20) -- (0,-1.65);
      \only<1>{
        \node[magic] at (0,-1.95)
          {$\mathsf{TEnc}(\ket{\mathrm{CCZ}})\;\text{or}\;\mathsf{rej}$};
      }
      \only<2-3>{
        \node[font=\large\bfseries,text=CVQCInk] at (0,-1.95) {?};
      }
    \end{scope}
  \end{tikzpicture}

  \note<1>{
    \footnotesize
    Existing authenticated-distillation approaches in quantum multiparty
    computation start with correctly authenticated states, though their
    plaintexts may be noisy. The red CCZ inside a black TEnc marks that
    distinction: the authentication is valid even for a bad plaintext.
    Authentication then protects the prescribed distillation operations,
    giving a modular analysis under the required joint input-quality
    promise. CCZ is a schematic magic-state label here; the complete
    construction first distills the other required authenticated magic
    states and gadgets. Citation: Dulek, Grilo, Jeffery, Majenz, and
    Schaffner, Secure Multi-party Quantum Computation with a Dishonest
    Majority, EUROCRYPT 2020. The manuscript cites this authenticated
    distillation strategy in its Technical Overview and Introduction.
  }
  \note<2>{
    \footnotesize
    Our MBQC preparation gives the factorizable output characterization,
    including a ciphertext-level attack outside the pad. It does not simply
    hand us an ideal QAS encoding of a noisy logical state. Cross out the
    old label and place a fully red TEnc expression above it to indicate
    that the authentication codeword itself may have been altered.
    This does not claim that every
    ciphertext is malformed or that no correct copies remain. The
    preparation and testing analysis still establishes the required
    structure and bounds on bad copies. The question mark records the
    distillation analysis that remains to be proved; it is not yet a
    clean-state-or-reject guarantee for arbitrary malformed inputs.
  }
  \note<3>{
    \footnotesize
    The authentication interface is not initially guaranteed to be clean,
    so we cannot invoke the usual modular argument to rule out cheating
    during distillation. The red process label marks this untrusted
    execution. This is an a priori challenge, not an impossibility claim:
    our analysis carries the outside attack and uses factorizability and
    iterative twirling to control the effect of further deviations. The
    eventual guarantee separates remaining ciphertext-level attacks from
    the plaintext faults that distillation removes. Next, show the basic
    four-to-one circuit on zero ancillas. Source: Technical Overview,
    Security of the Distillation Sub-Protocol, and the distillation section.
  }
\end{frame}
